\section{問題}

\subsection{2018}
\tbox{
\begin{prob}
三角形$T$の一つの辺の長さは平方数で, 残りの辺の長さは素数である. また, $T$の面積は整数で, 外接円の著系の長さは素数である. $T$の各辺の長さを求めよ.
\end{prob}
}{title=2018年3月号「読者と作るページ」掲載問題}

\tbox{
\begin{prob}
座標平面上の原点Oを中心とした半径1の円周を$C_1$, 曲線$y=kx^3$を$C_2$とする. ただし, $k>0$とする. $C_1$と$C_2$は第一象限において点Aで交わり, $C_1,C_2$の点Aにおける接線をそれぞれ$l,m$とする. $l,m$のなす角が最小となるように$k$をとるとき, 第一象限において, $C_1,C_2$,$y$軸で囲まれる領域の面積を求めよ.
\end{prob}
}{title=2018年3月号「読者と作るページ」掲載問題}

\tbox{
\begin{prob}
$xy$平面上の放物線$G$:$y=mx^2 + px + q^n$を考える. $G$と$x$軸は2つの交点をもつものとする. 各交点における$G$の接線2本と$x$軸で囲まれる部分の面積を$S$とするとき,
\begin{enumerate}
    \item $S$を$m,n,p,q$で表せ.
    \item $m$を整数, $n$を自然数, $p,q$を素数とする. $S$が自然数となるとき,
    \begin{enumerate}
        \item $p$を求めよ.
        \item $S$を求めよ.
    \end{enumerate}
\end{enumerate}
\end{prob}
}{title=2018年学コン9月号3番}
%----2019----

\newpage
\subsection{2019}
\tbox{
\begin{prob}
 $n$を正の整数, $e$を自然対数の底とする。
 \begin{enumerate}
  \item 関数$\disp f_n(x)=1+x+\frac{x^2}{2!}+\frac{x^3}{3!}+\dots +\frac{x^n}{n!}$ とするとき、$x>0$において、$0<e^x-f_n(x)<\dfrac{x^{n+1}e^x}{(n+1)!}$であることを示せ。
  \item $\disp \lim_{n\to\infty} \cos\frac{2\pi e n!}{3}=-\frac{1}{2}$ を示せ。
 \end{enumerate}
\end{prob}
}{title=2019年6月号学コン5番}

\tbox{
\begin{prob}
$20^m\div (1-9^n)$が整数となるような自然数の組$(m,n)$を求めよ.
\end{prob}
}{title=2019年9月号学コン5番}

\tbox{
\begin{prob}
$xyz$空間の点${\rm P}(\sin{\theta}, \sin^2{\theta}, 0)$と${\rm Q}(0,-\cos^2{\theta}, \cos{\theta})$について考える. 実数$\theta$が$0\leq \theta\leq \dfrac{\pi}{2}$の範囲で動くとき, 線分${\rm PQ}$の通貨する領域を$S$とする. $S$を$y$軸の周りに一回転させて得られる立体の体積を求めよ.
\end{prob}
}{title=2019年11月号学コン6番}



\clearpage
\subsection{2020}
%---2020---

\tbox{
\begin{prob}%Accepted
四面体OABCは, $\mr{OA}=1, \mr{OB}=\mr{OC}=2$を満たし, 面$\mr{ABC}$が正三角形であるとする。
\begin{enumerate}
\item 正三角形$\mr{ABC}$の一辺の長さ$x$の取りうる値の範囲を求めよ。
\item 四面体OABCの体積の最大値を求めよ。
\end{enumerate}
\end{prob}
}{title=2020年10月号学コン5番}


\tbox{
\begin{prob}
\ben
\item $x^2 + y^2 = 3z^2$の整数解は$(x,y,z) = (0,0,0)$ に限ることを証明せよ。
\item $xy$平面において, $x$座標と$y$座標がともに有理数であるような点のことを「有理点」と呼ぶこととする。$\theta$を$0< \theta \leq \dfrac{\pi}{2} $を満たす定数とし, 次の条件を満たす$xy$平面上の三角形$\mr{ABC}$を考える:
\begin{center}
(\tf{条件})\qquad $\mr{A}$は有理点であり, $\mr{AB}$の長さは有理数である。さらに, $\angle{\mr{ACB}} = \theta$である。
\end{center}
三角形$\mr{ABC}$の外心を$\mr{X}$とする。次の問に答えよ。
\ben
\item $\theta = \dfrac{\pi}{3}$のとき, $\mr{X}$が有理点でないことを証明せよ。
\item $\sin{\theta}= \dfrac{1}{\sqrt{2020}}$ のとき, $\mr{X}$が有理点となることはあるか. あるならば点$\mr{A},\mr{B},\mr{C}$の座標を挙げよ.
\een
\een
\end{prob}
}{title=2020年11月号学コン5番}

\tbox{
\begin{prob}
$f(x)=\dfrac{1}{1+x-x^3}$とおく. 素数$p$, 正の整数$k$, および$p$で割れない整数$m$の組$(p,k,m)$であって, $f\parena{\dfrac{m}{p^k}}$が整数となるものをすべて求めよ.
\end{prob}
}{title=2020年12月号宿題}

\clearpage
\subsection{2021}
%---2021-----
\tbox{
\begin{prob}
$\dfrac{n!}{m^{n-1}}$が整数であるような2以上の整数$m,n$の組を求めよ.
\end{prob}
}{title=2021年6月号学コン6番}


\clearpage
\subsection{2022}

\tbox{
\begin{prob}
原点を$\mr{O}$とする座標平面上の点$\mr{A}_n$ ($n=0,1,2,\dots$)を以下の操作で定める.
\[\mr{A}_0 = (1,0),\quad \mr{A}_{1} = (\cos{\dfrac{2\pi}{N}}, \sin{\dfrac{2\pi}{N}})\quad (N\text{は正の整数})  \]
\[\text{$n\geq 2$に対し, $\mr{A}_{n}$は直線$\mr{OA}_{n-2}$に関して点$\mr{A}_{n-1}$と対称な点.}\]
\begin{enumerate}
\item $\mr{A}_n$の$x$座標$x_n$を$n$を用いて表せ.
\item 数列$x_n$ ($n=0,1,2,\dots$)が収束するような正の整数$N$の条件を求めよ.
\end{enumerate}
\end{prob}
}{title=2022年学コン11月号}


\clearpage
\subsection{2023}
%2023
\tbox{
\begin{prob}
表と裏が等確立で出るコインを$n$回投げて, 点${\rm P}_n(x_n,y_n,z_n)$を次のように定める.
\begin{align*}
&(x_0,y_0,z_0) = (1,-3,1) \\
&\text{$k$回目で表が出たとき}(x_k,y_k,z_k) = (-x_{k-1}, -y_{k-1} - z_{k-1} , x_{k-1}+y_{k-1} + z_{k-1}) \\
&\text{$k$回目で表が出たとき}(x_k,y_k,z_k) = (z_{k-1}, y_{k-1}, x_{k-1})
\end{align*}
\begin{enumerate}
\item ${\rm P}_n$が$(1,-1,0)$である確率を$a_n$, $(-1,1,0)$である確率を$b_n$として$a_n+b_n$を$n$で表せ.
\item ${\rm P}_n$が$(0,0,0)$である確率を$n$で表せ.
\end{enumerate}
\end{prob}
}{title=2023年1月号学コン6番}


\tbox{
\begin{prob}
    $f(x)$は1次以上の実数係数の多項式で, すべての正の整数$n$に対して$f(n)$は正の整数であるとする. このとき, $f(n)$を10進法で表したときの桁として$0$が現れるような正の整数$n$が存在することを証明せよ.
\end{prob}
}{title=2023年3月号宿題}


\tbox{
\begin{prob}
$a,b$が実数全体を動くとき, $xy$平面上の点$\mr{P}(a+b+ab, (a+b)ab)$が動く範囲を求めよ.
\end{prob}
}{title=2023年5月号学コン6番}


\tbox{
\begin{prob}
$0 \leq \theta < \dfrac{\pi}{2}$を満たす実数$\theta$が次の条件をみたすならば, $\theta=0$であることを示せ.
\begin{center}
$\tan{\dfrac{\theta}{2^n}}$が有理数となるような非負整数$n$が無数に存在する.
\end{center}
\end{prob}
}{title=2023年7月号宿題}


\tbox{
\begin{prob} %Accepted
実部と虚部がともに整数であるような複素数全体の集合を$A$で表す. また, $a=1+i$とおく(ただし$i=\sqrt{-1}$). 以下の問に答えよ.
\begin{enumerate}
\item $A$の元$z_1,z_2$に対し, $\dfrac{z_1z_2}{a}\in A$であるとする. このとき, $\dfrac{z_1}{a}$と$\dfrac{z_2}{a}$のいずれか一方は$A$に属すことを証明せよ.
\item $x^3 + ay^3 + a^2 z^3 + axyz = 0 \, \dots\textcircled{1}$を満たす$A$の要素の組$(x,y,z)$を考える.
\begin{enumerate}
\item $(x,y,z)=(\alpha,\beta,\gamma)$のとき, $(x,y,z)=(\dfrac{\alpha}{a},\dfrac{\beta}{a},\dfrac{\gamma}{a})$も\textcircled{1}を満たす$A$の要素であることを示せ.
\item 組$(x,y,z)$をすべて求めよ.

\end{enumerate}

\end{enumerate}
\end{prob}
}{title=2023年11月号学コン6番}

\tbox{
\begin{prob}
%\begin{enumerate}
%\item 同一平面上に正六角形と, 相異なる3直線がある. 正六角形の6頂点は3直線のいずれかに含まれており, 3直線はある1点を共有しているとする. このとき, 3直線は正六角形の最長の対角線3本を延長して得られる直線であることを示せ.
%\item
$xyz$空間において$y^3 + z(x^2 + z^2)=0$で表される曲面を$S$とする. 1辺の長さが1の正六角形で, 2頂点が$x$軸上に, 4頂点が$S$上にあるものがただ一つ存在することを証明せよ.
%\end{enumerate}
\end{prob}
}{title=2023年11月号宿題}

\clearpage
\subsection{2024}
%-----------------2024--------------------
\tbox{
\begin{prob} %Accepted
$\sqrt[abc]{a! + b! + c!}$が正の整数であるような正の整数の組$(a,b,c)$をすべて求めよ.
\end{prob}
}{title=2024年2月号宿題}

\tbox{
\begin{prob}
$xyz$空間において連立方程式
\[
y^2 = x^2(x+1), \quad y=zx
\]
で表される集合を$C$とする. 原点を中心とする半径$r>0$の球面を$S$とするとき, $C$と$S$の共有点はいくつあるか. $r$で場合分けして答えよ.
\end{prob}
}{title=2024年4月号学コン4番}



% \tbox{
% \begin{prob}
% $\varphi = \dfrac{1+\sqrt{5}}{2}$とする. 極限$\disp\lim_{n\to \infty} \sin{(2\pi r \varphi^n)}$が存在する有理数$r$をすべて求めよ.
% \end{prob}
% }{title=2024年6月号宿題}

% \tbox{
% \begin{prob} %Accepted
% 辺の長さがすべて1の正三角柱$\mr{ABC}-\mr{DEF}$がある. 端点を含む線分$\mr{AE}, \mr{BF}, \mr{CD}$上にそれぞれ点$\mr{P}, \mr{Q}, \mr{R}$をとり, 三角形$\mr{PQR}$の重心を$\mr{G}$とする. $\mr{P}, \mr{Q}, \mr{R}$を自由に動かしたときに$\mr{G}$が通過する領域を$K$とおく. $0\leq t\leq 1$を満たす実数$t$に対し, 面${\rm ABC}$を底面として高さ$t$の所にある$K$の点のなす領域について, その面積を$S(t)$とおく. $t\leq\frac{1}{3}$であるとき, $S(t) + S(t+\frac{1}{3}) + S(t+\frac{2}{3})$をもとめよ.
% \end{prob}
% }{title=2024年6月号学コン5番}

